% GENERATED FILE. Edit canonical lessons, then run npm run generate:books.
% canonical-source-hash: fnv1a64:e6c5dbcc60f704d0
% canonical-lessons: ZH-C533-xiangwang, ZH-C533-xiwang, ZH-W533-xi-hope, ZH-C533-kewang, ZH-C533-juewang, ZH-C533-xinxiangshicheng

\chapter{I Hope So}
\label{ch:zh-i-hope-so}
\hlchaptermodality{533}

\begin{chapteropening}
\textbf{By the end of this chapter:} \emph{I can say what I hope, long for and yearn for, say when hope runs out, and wish someone well.}

Complete the last lesson of chapter 533: I can say what I hope, long for and yearn for, say when hope runs out, and wish someone well.
\end{chapteropening}

\section[xiàngwǎng]{xiàngwǎng — to yearn for}

\label{lesson:ZH-C533-xiangwang}



\begin{quote}
Before the new one: say the Mandarin for as one would wish it, then the Mandarin for to count on someone.
\end{quote}

\subsection*{You'll want to know: xiàngwǎng}
\emph{xiàngwǎng} — \textquotedblleft{}to yearn for\textquotedblright{}.

\begin{quote}
\emph{rúguǒ yǒu jīhuì wǒ jiù qù wǒ xiàngwǎng de dìfang} — If I get the chance, I'll go to the place I've always longed for.
\end{quote}

Its characters come later; for now it is a word for the ear.

\begin{cousinweb}
\emph{xiàngwǎng} is built from \emph{xiàng} 'towards' and \zh{往} 'to go', both inherited from Middle Chinese, \emph{xjangH} and \emph{hjwangX} (Guangyun, in Baxter's transcription), which gives the first as the same word as its usual spelling and glosses the second 'to go to'.
\end{cousinweb}

\begin{grammarlens}[title={What you've built}]
One more word for saying what would happen.
\end{grammarlens}

\subsection*{Your turn}
\begin{itemize}
\raggedright
  \item \emph{Say it:} \emph{xiàngwǎng}
  \item \emph{Say it:} \emph{xiàngwǎng}, once more
  \item \emph{Say it:} \emph{rúguǒ yǒu jīhuì wǒ jiù qù wǒ xiàngwǎng de dìfang} — If I get the chance, I'll go to the place I've always longed for.
  \item \emph{Recall:} say the Mandarin for as one would wish it, then the Mandarin for to count on someone, then say \emph{xiàngwǎng} again
\end{itemize}

\begin{tcolorbox}[breakable,colback=teal!4,colframe=teal!35!black,title={Before you move on}]
What does xiàngwǎng mean? (\textquotedblleft{}To yearn for\textquotedblright{}.) Say it once more.
\end{tcolorbox}
\section[\texorpdfstring{\zh{希望} (xīwàng)}{希望 (xīwàng)}]{\zh{希望} (xīwàng) — to hope}

\label{lesson:ZH-C533-xiwang}



\begin{quote}
Before the new one: say the Mandarin for to count on someone, then the Mandarin for to yearn for.
\end{quote}

\subsection*{You'll want to know: \zh{希望}}
\textbf{\zh{希望}} — \emph{xīwàng} — \textquotedblleft{}to hope\textquotedblright{}.

\begin{quote}
\textbf{\zh{我希望明天不下雨}} \emph{wǒ xīwàng míngtiān bú xiàyǔ} — I hope it doesn't rain tomorrow.
\end{quote}

The new character, in pieces: \zh{希} — seven strokes in three pieces: a small cross of two strokes on top; a long bar with a long stroke falling to the left through it; and at the foot three strokes, a short vertical, a bar that turns down and hooks, and a long vertical through the middle. The next lesson puts it in your hand.

\begin{cousinweb}
\textbf{\zh{希望}} (\emph{xīwàng}) is built from \zh{希} and \zh{望} 'to look out for', both inherited from Middle Chinese, \emph{xj+j} and \emph{mjangH} (Guangyun, in Baxter's transcription). The dictionary glosses the first 'to stop', 'to scatter' and, after an older dictionary, 'rare', and also with the second: 'to hope'.
\end{cousinweb}

\begin{grammarlens}[title={What you've built}]
One more word for saying what would happen.
\end{grammarlens}

\subsection*{Your turn}
\begin{itemize}
\raggedright
  \item \emph{Say it:} \emph{xīwàng}
  \item \emph{Say it:} \emph{xīwàng}, once more
  \item \emph{Say it:} \emph{wǒ xīwàng míngtiān bú xiàyǔ} — I hope it doesn't rain tomorrow.
  \item \emph{Recall:} say the Mandarin for to count on someone, then the Mandarin for to yearn for, then say \emph{xīwàng} again
\end{itemize}

\begin{tcolorbox}[breakable,colback=teal!4,colframe=teal!35!black,title={Before you move on}]
What does \zh{希望} mean? (\textquotedblleft{}To hope\textquotedblright{}.) Say it once more.
\end{tcolorbox}
\section[\texorpdfstring{\zh{希} (xī)}{希 (xī)}]{\zh{希} — a cross, a bar and a long stroke, then three at the foot}

\label{lesson:ZH-W533-xi-hope}



\begin{quote}
The word first, then the character inside it.

\begin{itemize}
\raggedright
  \item \emph{Recall:} say \emph{xīwàng} — \textbf{R1}, one lesson back, before it decays
\end{itemize}

\begin{quote}
\zh{望} \zh{有}
\end{quote}

\emph{Read it:} each one aloud, with nothing said first — \emph{wàng}, \emph{yǒu}
\end{quote}

\subsection*{Script — a cross, a bar and a long stroke, then three at the foot}
\hlblockfigure{\detokenize{figures/ZH-W533-xi-hope-filmstrip.pdf}}{How \zh{希} is written, stroke by stroke}

\begin{quote}
\zh{希}
\end{quote}

\textbf{Three pieces, top to bottom.} Count the pieces, not the strokes.

\textbf{The top piece, two strokes}: a short stroke falling to the left, then a short stroke falling to the right across it, a small cross.

\textbf{The middle piece, two strokes}: a long bar, then a long stroke falling to the left from above the bar down through it.

\textbf{The bottom piece, three strokes}: a short vertical on the left, a bar that turns down the right side and hooks to the left, then a long vertical down through the middle, past the foot.

2 + 2 + 3 = seven.

\subsection*{Your turn}
\begin{itemize}
\raggedright
  \item \emph{Write it:} \zh{希}, all seven, with the model in front of you
  \item \emph{Cover:} write it again from memory, and count the pieces, not the strokes
\end{itemize}

\begin{tcolorbox}[breakable,colback=teal!4,colframe=teal!35!black,title={Before you move on}]
How many strokes does \zh{希} have? (\textbf{Seven}.) What is the last stroke? (\textbf{A long vertical through the middle of the foot}.)

Source: \href{https://raw.githubusercontent.com/chanind/hanzi-writer-data/68d10a4b21150cae5e1ebbd223eed289cf32d90c/data/\%E5\%B8\%8C.json}{Hanzi Writer Data \zh{希}, pinned stroke paths and medians}.
\end{tcolorbox}
\section[kěwàng]{kěwàng — to thirst for; to crave}

\label{lesson:ZH-C533-kewang}



\begin{quote}
Before the new one: say the Mandarin for to yearn for, then the Mandarin for to hope that.
\end{quote}

\subsection*{You'll want to know: kěwàng}
\emph{kěwàng} — \textquotedblleft{}to thirst for; to crave\textquotedblright{}.

\begin{quote}
\emph{tā hěn kěwàng dāng yīshēng} — He longs to be a doctor.
\end{quote}

Its characters come later; for now it is a word for the ear.

\begin{cousinweb}
\emph{kěwàng} is built from \emph{kě} 'thirsty' and \zh{望} 'to look out for', both inherited from Middle Chinese, \emph{khat} and \emph{mjangH} (Guangyun, in Baxter's transcription), which glosses the first 'hunger and thirst'.
\end{cousinweb}

\begin{grammarlens}[title={What you've built}]
One more word for saying what would happen.
\end{grammarlens}

\subsection*{Your turn}
\begin{itemize}
\raggedright
  \item \emph{Say it:} \emph{kěwàng}
  \item \emph{Say it:} \emph{kěwàng}, once more
  \item \emph{Say it:} \emph{tā hěn kěwàng dāng yīshēng} — He longs to be a doctor.
  \item \emph{Recall:} say the Mandarin for to yearn for, then the Mandarin for to hope that, then say \emph{kěwàng} again
\end{itemize}

\begin{tcolorbox}[breakable,colback=teal!4,colframe=teal!35!black,title={Before you move on}]
What does kěwàng mean? (\textquotedblleft{}To thirst for; to crave\textquotedblright{}.) Say it once more.
\end{tcolorbox}
\section[juéwàng]{juéwàng — to give up all hope; to despair}

\label{lesson:ZH-C533-juewang}



\begin{quote}
Before the new one: say the Mandarin for to hope that, then the Mandarin for to thirst for.
\end{quote}

\subsection*{You'll want to know: juéwàng}
\emph{juéwàng} — \textquotedblleft{}to give up all hope; to despair\textquotedblright{}.

\begin{quote}
\emph{jíshǐ shībài le yě búyào juéwàng} — Even if you fail, don't give up hope.
\end{quote}

Its characters come later; for now it is a word for the ear.

\begin{cousinweb}
\emph{juéwàng} is built from \emph{jué} 'to cut off' and \zh{望} 'a hope', both inherited from Middle Chinese, \emph{dzjwet} and \emph{mjangH} (Guangyun, in Baxter's transcription), which glosses the first 'to break off'.
\end{cousinweb}

\begin{grammarlens}[title={What you've built}]
One more word for saying what would happen.
\end{grammarlens}

\subsection*{Your turn}
\begin{itemize}
\raggedright
  \item \emph{Say it:} \emph{juéwàng}
  \item \emph{Say it:} \emph{juéwàng}, once more
  \item \emph{Say it:} \emph{jíshǐ shībài le yě búyào juéwàng} — Even if you fail, don't give up hope.
  \item \emph{Recall:} say the Mandarin for to hope that, then the Mandarin for to thirst for, then say \emph{juéwàng} again
\end{itemize}

\begin{tcolorbox}[breakable,colback=teal!4,colframe=teal!35!black,title={Before you move on}]
What does juéwàng mean? (\textquotedblleft{}To give up all hope; to despair\textquotedblright{}.) Say it once more.
\end{tcolorbox}
\section[\texorpdfstring{\zh{心想事成} (xīnxiǎng shìchéng)}{心想事成 (xīnxiǎng shìchéng)}]{\zh{心想事成} (xīnxiǎng shìchéng) — may your wishes all come true}

\label{lesson:ZH-C533-xinxiangshicheng}



\begin{quote}
Before the new one: say the Mandarin for to thirst for, then the Mandarin for to give up all hope.
\end{quote}

\subsection*{You'll want to know: \zh{心想事成}}
\textbf{\zh{心想事成}} — \emph{xīnxiǎng shìchéng} — \textquotedblleft{}may your wishes all come true\textquotedblright{}.

\begin{quote}
\textbf{\zh{祝你们好运心想事成}} \emph{zhù nǐmen hǎoyùn xīnxiǎng shìchéng} — Good luck to you all, and may all your wishes come true.
\end{quote}

\begin{cousinweb}
\textbf{\zh{心想事成}} (\emph{xīnxiǎng shìchéng}) is built from \zh{心} 'the heart', \zh{想} 'to wish', \zh{事} 'a matter' and \zh{成} 'done', all inherited from Middle Chinese, \emph{sim}, \emph{sjangX}, \emph{dzriH} and \emph{dzyeng} (Guangyun, in Baxter's transcription): what the heart thinks of, may it be done.
\end{cousinweb}

\begin{grammarlens}[title={What you've built}]
One more word for saying what would happen.
\end{grammarlens}

\subsection*{Your turn}
\begin{itemize}
\raggedright
  \item \emph{Say it:} \emph{xīnxiǎng shìchéng}
  \item \emph{Say it:} \emph{xīnxiǎng shìchéng}, once more
  \item \emph{Say it:} \emph{zhù nǐmen hǎoyùn xīnxiǎng shìchéng} — Good luck to you all, and may all your wishes come true.
  \item \emph{Recall:} say the Mandarin for to thirst for, then the Mandarin for to give up all hope, then say \emph{xīnxiǎng shìchéng} again
\end{itemize}

\begin{tcolorbox}[breakable,colback=teal!4,colframe=teal!35!black,title={Before you move on}]
What does \zh{心想事成} mean? (\textquotedblleft{}May your wishes all come true\textquotedblright{}.) Say it once more.
\end{tcolorbox}
